\section*{Introduction}
\addcontentsline{toc}{section}{Introduction}

Ce document est une annexe au rapport de projet de la fusée bi-étage SP-02, développée par l'association \acrshort{IPSA} AéroIPSA
au cours de l'année scolaire 2025-2026 \cite{rapportSP02}. Il n'a pas vocation à présenter le projet dans son ensemble :
l'architecture générale de la fusée, le dimensionnement de ses sous-systèmes et le déroulement de la mission y sont supposés
connus. La lecture préalable du rapport principal est donc nécessaire à la bonne compréhension du présent document, qui n'en
développe qu'un point particulier.\\

Pour situer brièvement le sujet traité ici, SP-02 est une fusée expérimentale à deux étages dont la séparation intervient après
l'extinction du propulseur du premier étage. Cette séparation repose sur un bloc aérofreins, embarqué dans le premier étage et
piloté par le séquenceur, dont le déploiement provoque une augmentation contrôlée de la traînée de l'étage inférieur. Le
différentiel de décélération ainsi créé entre les deux corps assure leur dégagement mutuel. L'efficacité du mécanisme de
séparation dépend donc directement des efforts aérodynamiques générés par les aérofreins, et c'est à la caractérisation de ces
efforts qu'est consacrée cette étude.\\

Une campagne d'essais en soufflerie ne pouvant être menée dans le cadre du projet (contraintes logistiques et de calendrier), la
simulation numérique constitue le seul moyen d'anticiper le comportement du système en conditions de vol. L'étude aérodynamique
par \acrfull{CFD} a ainsi été conduite en parallèle de la conception mécanique du bloc aérofreins, dans le but de valider
théoriquement le concept avant sa mise en \oe uvre.\\

L'approche \acrshort{CFD} présente pour ce projet plusieurs intérêts. Elle donne accès à des grandeurs difficilement mesurables
sans moyens d'essai dédiés, en particulier les champs de pression et de vitesse ainsi que les efforts aérodynamiques localisés sur
les aérofreins. Elle permet de comparer quantitativement plusieurs configurations d'ouverture et d'en déduire une loi reliant le
taux de déploiement à la traînée du premier étage. Elle autorise enfin l'analyse de la structure de l'écoulement autour de la
fusée, notamment le développement et le décollement de la couche limite, phénomène moteur dans la génération de traînée.\\

L'étude se limite volontairement à l'aérodynamique du système. La tenue mécanique des pièces du bloc aérofreins, la cinématique de
leur déploiement et le pilotage du servomoteur associé relèvent d'autres parties du projet et ne sont pas abordés ici. De même,
seules des positions d'ouverture figées sont considérées, la phase transitoire de déploiement sortant du périmètre de ce travail.\\

Le document s'organise en deux parties. La première décrit la méthodologie de simulation retenue : simplification de la maquette
numérique, hypothèses de modélisation physique, construction du domaine de calcul et stratégie de maillage. La seconde présente et
analyse les résultats obtenus, à travers l'évolution du coefficient de traînée, la répartition de pression sur les aérofreins et
la structure des champs de vitesse.

\newpage
